\clearpage
\fancyhead[L]{\small HARDING. CELESTIAL CHARTS.}
\hypertarget{native-598-lower}{}
\begin{center}\small Göttingische gelehrte Anzeigen. No. 5. Pages 41..45. 9 January 1817.\\[3pt]\rule{.72\linewidth}{.3pt}\end{center}

A new instalment of our Professor \textsc{Harding}’s celestial charts, the latest instalments of which we announced in No. 158 of this journal for 1812 [p. 541 of this volume], has now appeared, again on commission through the Vandenhoeck-Ruprecht bookshop. This \emph{fifth} instalment comprises four sheets, namely Nos. VI, X, XIX and XXV. To make it easier to see how much of this useful work has now been completed and how much still remains, we confine ourselves here to a brief reminder of its origins and plan. The discovery of the planets Ceres and Pallas, which never appear larger than fixed stars of the eighth magnitude and often appear much smaller still, first made far more detailed star charts than we had possessed up to that time indispensable for the parts of the sky in which those interesting celestial bodies can appear. To avoid needing separate charts for them every year, it was desirable to have an atlas devoted to the zodiac of these planets, and the author decided to prepare one. A moderate number of sheets (10 or 11) would have sufficed for this. But while the author was already pursuing this work, and partly as a result of this very work, two other new planets were also discovered, making it necessary to extend the plan. The many benefits that detailed star \SourceLeaf{599}charts also afford in other respects prompted the author to extend his atlas, some sheets of which had already been completed in the meantime, to the whole northern sky and to the part of the southern sky that can still be observed well from our parts of Europe and for which the materials are still available, thus to a little more than $30^\circ$ south declination. Nine sheets, each spanning $40^\circ$ in right ascension, were therefore to cover this southern part of the sky, and nine sheets, each with a similar span, were likewise to cover the part of the northern sky from the equator to $32^\circ$ north declination. These 18 sheets are now all complete, with the addition in the present instalment of sheets VI and X, which had hitherto been missing. However, sheet No. I was drawn up at the outset, when the limited plan was still in force. It extends from nearly $22^\circ$ south declination to almost $13^\circ$ north declination, and its northern part is also included on the new sheet No. X. Sheet No. I will in future be withdrawn and replaced by another extending to $32^\circ$ south declination. The author will shortly supply this replacement to the owners of the atlas free of charge.

The region of the sky covered by these 18 sheets is of greatest interest to astronomers because the planets and the Moon never go beyond it. The author therefore rightly believed that he should complete it first. Yet an equally detailed survey of the northern sky from $32^\circ$ declination to the pole would have to be regarded as essential if only because comets so often pass through it. Many comets are visible \emph{only} through a telescope, and even the larger ones can still be followed for a time with telescopes once their increasing distance has already taken them beyond the reach of the naked eye. On such occasions, for lack of engraved star charts, the observer always has to undertake the time-consuming work of drawing a chart of the region through which the comet passes. Sheets XIX and XXV of the present atlas are a first step towards meeting this need. The plan followed in dividing up the northern sky is as follows. Since the degrees along the parallels become smaller towards the pole, a sheet of the same size as the first 18 could cover more in right ascension. Eight sheets will follow first, each extending a little over $45^\circ$ in right ascension at its lower edge and from $30^\circ$ to nearly $65^\circ$ in declination in the middle. The remaining part of the sky, from $64^\circ$ declination to the pole, can then be covered on two sheets at a scale only slightly reduced. The corners, which would otherwise repeat regions of the sky already shown on the sheets, are to be used to depict interesting groups of stars such as the Pleiades, the Manger and others. The entire Athas will thus be complete in 28 sheets, and it will hopefully be possible to supply shortly the eight sheets still missing, most of which have already been drawn.

For the first 18 sheets a more elaborate projection was unnecessary. The grid of plain squares made the charts easier to draw and also offers several advantages when they are used. Near the boundary the slight and scarcely perceptible departure from the true ratio of degrees of right ascension to degrees of declination was of no consequence at all. At larger declinations, however, this was no longer applicable. For the eight sheets XIX to XXVI the author therefore chose a projection proposed to him for this purpose by Councillor \textsc{Gauss}, in which the circles of declination are straight lines and the parallels are concentric circles. Their decreasing radii are determined with full precision so that degrees of right ascension and degrees of declination are everywhere in the correct proportion. In this representation the image is therefore perfectly similar to the sky in its smallest parts, which is the most essential condition, \SourceLeaf{600}Strictly speaking, the degrees of declination are admittedly not all the same size. But this circumstance, which is unavoidable by the nature of the problem, is of no consequence here, especially since the difference is so small that it is scarcely noticeable. At the 64th and 3oth degrees of declination, one degree of declination measures 13.07 millimetres, whereas in the middle, at the 47th degree, it measures 12.38 millimetres.

Since the division of the sky among the charts is determined solely by right ascensions and declinations, some constellations naturally appear in pieces. Sheet No.\ VI contains parts of the Virgin, the Balance, the Scorpion, the Serpent, the Wolf, the Centaur and the tail of Hydra, and all of the Hermit Bird. Sheet No.\ X contains parts of the Fishes, Andromeda, the Triangles, the Ram, the Whale and Pegasus. Sheet No.\ XIX contains parts of Cepheus, Cassiopeia, the Harvest Keeper, the Camelopard, Perseus, Andromeda, the Fishes, the Triangles and the Ram. Sheet No.\ XXV contains parts of the Dragon, Hercules, the Lyre, the Swan and Cepheus. The paper and engraving are as fine as those of the earlier instalments.
